\documentclass[11pt]{article}
\usepackage[T1]{fontenc}
\usepackage{lmodern,microtype,amsmath,amssymb,amsthm,mathtools}
\usepackage{booktabs,array,longtable,enumitem,xcolor}
\usepackage[a4paper,margin=25mm,top=25mm,bottom=25mm]{geometry}
\usepackage{fancyhdr}
\usepackage[colorlinks=true,linkcolor=blue!45!black,citecolor=teal!65!black,urlcolor=teal!65!black]{hyperref}
\definecolor{ink}{HTML}{183D4B}
\definecolor{muted}{HTML}{516772}
\setlength{\headheight}{15pt}
\setlength{\emergencystretch}{3em}
\setlength{\parskip}{3pt}
\setlist{itemsep=3pt,topsep=5pt}
\pagestyle{fancy}
\fancyhf{}
\fancyhead[L]{\small\color{ink}Beyond Dirichlet GRH}
\fancyhead[R]{\small\color{muted}Transfer analysis}
\fancyfoot[C]{\small\thepage}
\raggedbottom
\newtheorem{theorem}{Theorem}[section]
\newtheorem{lemma}[theorem]{Lemma}
\newtheorem{proposition}[theorem]{Proposition}
\newtheorem{corollary}[theorem]{Corollary}
\theoremstyle{definition}
\newtheorem{definition}[theorem]{Definition}
\newtheorem{remark}[theorem]{Remark}
\newcommand{\Q}{\mathbb Q}
\newcommand{\Z}{\mathbb Z}
\newcommand{\F}{\mathbb F}
\newcommand{\C}{\mathbb C}
\newcommand{\Cl}{\operatorname{Cl}}
\newcommand{\ord}{\operatorname{ord}}
\newcommand{\Hom}{\operatorname{Hom}}
\newcommand{\vp}{v_p}
\newcommand{\supp}{\operatorname{supp}}
\newcommand{\N}{\operatorname{N}}
\newcommand{\rank}{\operatorname{rank}}
\newcommand{\code}[1]{\texttt{\small\detokenize{#1}}}
\hypersetup{pdftitle={Beyond Dirichlet GRH: arithmetic twists and sparse equations},pdfauthor={Research prepared for Vladimir Reshetnikov}}
\begin{document}
\thispagestyle{plain}
\noindent{\small\sffamily\color{ink}RESEARCH REPORT \hfill 25 NOVEMBER 2026}
\vspace{15pt}
\begin{center}
{\LARGE\bfseries\color{ink} Beyond Dirichlet GRH}\\[7pt]
{\Large Arithmetic twists, sparse equations,\\[3pt]and the remaining transfer barriers}\\[9pt]
{\normalsize Prepared for Vladimir Reshetnikov}
\end{center}
\vspace{8pt}
\noindent\textbf{Abstract.}
We examine the arithmetic mechanism of the shorter Lean proof of Dirichlet GRH and its prospects for further mathematical applications. The source is pinned to revision \code{72fead102f663d5c7a407964a13dc0cf9d977207}; the comparison corpus is the 41 incoming ProveIt archives at revision \code{48efe71a726dddc5a505534edaa4c0fd190416b9}. We give an exact criterion for turning selected local torsion characters into global anticyclotomic characters, including prescribed splitting conditions, and quantify the possible loss in order and independent parameters. An exact example shows that a local cubic character at 13 over $\Q(\sqrt{-23})$ requires global order nine, while imposing splitting at 3 eliminates all nontrivial characters on that support. The least conductor norm after this condition is imposed is $151^2=22801$, within the stated class of split-supported characters.

We also strengthen an incoming confluent determinant criterion to deeper congruences and small primes, prove the sharp blockwise threshold, and give exact seven-term support results for every positive power of two explicit powered singular moduli. Finally, an abstract covolume calculation shows how a better bound for completing a partial basis could replace the upstream requirement of nearly complete spanning.

The new results are paper proofs with exact computational certificates. Two previous CM certificate suites were fully rerun; their five-term classification is prior work, not a new claim of this report. We do not establish nonabelian GRH, BSD, Vandiver's conjecture, or a new general $K$-group vanishing theorem. The main contribution is a set of precise arithmetic interfaces, concrete calculations, and a more selective next research target.
\clearpage
\begingroup
\small
\setlength{\parskip}{0pt}
\tableofcontents
\endgroup
\newpage
\section{Assessment and what is added here}

There are useful transfers beyond the usual consequences of RH and Dirichlet GRH. Their present strength varies substantially. Some of the incoming work already proves concrete statements using additional local arithmetic; some extracts reusable lemmas; and some gives a logically exact criterion for an open problem while leaving its difficult arithmetic estimate unproved. These categories must remain separate.

The most substantial completed result that we re-examined is the incoming classification of rational equations with at most five nonzero terms for singular moduli. The proof combines local root counts, complex conjugate separation, and finite tests that remain valid for unbounded exponents. Its certificate suites reproduced successfully. This is more than substituting GRH into a familiar corollary: the strongest five-term classification in that archive is independent of GRH.\cite{prior_sparse}

This report adds the following.

\begin{enumerate}
\item \textbf{An exact arithmetic compatibility test.} A local torsion tuple extends with the same global exponent and specified triviality conditions precisely when it annihilates a particular intersection subgroup. The quotient class group explains the possible increase in order. The CM version also retains the root-of-unity relation which disappears only at suitable odd primes.
\item \textbf{A complete small conductor calculation.} For $\Q(\sqrt{-23})$, the minimum conductor norm of a nontrivial finite $3$-power anticyclotomic character, trivial on decomposition groups above $3$ and ramified only at split pairs, is $22801$. Exactly two characters attain it. The example distinguishes a local character problem from the more restrictive character problem actually used in the Lean source.
\item \textbf{A stronger unbounded-exponent certificate.} In a finite free $\Z_{(p)}$-algebra, a scalar congruence to depth $s$ permits a confluent test whenever each residue block has fewer than $p^s$ exponents. This includes $p=2$. We prove the threshold is sharp in this generality and obtain exact support seven for the families in Section~\ref{sec:confluent}.
\item \textbf{A quantitative alternative to almost-full spanning.} The determinant argument only charges missing dimensions by their excess completion cost. A cheaper completion could compensate for a much larger rank deficit. This identifies a distinct arithmetic lemma worth seeking for a new target.
\end{enumerate}

The finite-group and class-field arguments use established mathematics. We claim a useful synthesis and explicitly verified instances, not priority for the underlying theory. The strengthened confluent argument is derived here from the incoming proof; a literature search is not a proof of global novelty.

\subsection{The current incoming corpus matters}

The pinned directory contains 41 ZIP archives. Earlier summaries that stop at the unresolved discriminant $-131$, or describe all transfers as unfinished nonabelian criteria, are outdated relative to this snapshot. In particular, \code{grh_sparse_singular_moduli.zip} closes the unbounded-exponent $-131$ case and the four remaining sextic discriminants in the global five-term problem.\cite{prior_sparse}

For an algebraic number $\alpha$, define
\[
 s(\alpha)=\min_{\substack{0\ne P\in\Q[X]\\P(\alpha)=0}}
             \#\supp(P).
\]
For a singular modulus of discriminant $\Delta$ this depends only on $\Delta$, and is denoted $s(\Delta)$. The incoming theorem states
\begin{equation}\label{eq:prior-classification}
 s(\Delta)\le5\quad\Longleftrightarrow\quad h(\Delta)\le4,
\end{equation}
with $s(\Delta)=5$ at class number four and $s(\Delta)=6$ at class number five. The exceptional singular modulus $0$ has support one and is consistent with the displayed equivalence.

Our audit reran the complete current CM suite and its older companion. It reproduced the exact polynomial checks, the 25 cubic, 84 quartic and 29 quintic order censuses, both implementations of the last quintic sieve, the global reduction through 36,608 candidates, and the unbounded-exponent tests for $-87,-348,-432,-675$. We also read the local and global reductions. These are paper proofs and executable finite certificates, not Lean formalizations. The supporting classical class-number and CM theorems are cited in the original reports; we did not reprove those external inputs.

Other archives contain claims about progression-free densities, prime races, class groups, CM forms, and arithmetic classifications. We inspected their inventory and selected source sections to avoid duplicate work. We do not certify every mathematical claim in all 41 archives. This report's independent replay claim is limited to the two identified CM packages.

\subsection{Answers to the research question}

\begin{center}
\small
\begin{tabular}{p{.24\linewidth}p{.34\linewidth}p{.32\linewidth}}
\toprule
Direction & Present concrete achievement & What remains\\
\midrule
Sparse CM equations & Previous global five-term theorem; new deeper local certificates and powered sextic families & Six terms for all degree-six singular moduli, then a uniform larger threshold\\
Auxiliary characters & Exact local/global obstruction, order bound, parameter bound, minimal example & Compatible Selmer witnesses and quantitative special Frobenius selection\\
Nonabelian zero theorems & Reusable analytic and covolume interfaces; earlier explicit splitting-statistic criteria & Actual arithmetic lattices, gains, completion and denominator estimates\\
$K$-theory / Vandiver & Source odd-cocycle bounds and earlier finite candidate calculations & Comparison with conventional cohomology and complete obstruction coverage\\
BSD & A clear eligibility test for attempted substitutions & New input overcoming the conjugate-dual and scalar-inertia barriers\\
\bottomrule
\end{tabular}
\end{center}

The strongest recommendation is to pursue the next support threshold for sparse CM equations while formalizing the new local certificate. The nonabelian zero problem remains the larger conceptual opportunity, but it should advance through an actual arithmetic construction rather than another equivalent prime-sum criterion.

\section{What the repository proves and how the method works}

\subsection{Endpoint and verification boundary}

The source's public statement file imports Mathlib and states zeros of the Riemann zeta function and complex Dirichlet $L$-functions. The five exports are:
\begin{itemize}
\item \code{riemannHypothesis};
\item \code{riemannZeta_eq_zero_re};
\item \code{LFunction_ne_zero_of_half_lt_re};
\item \code{LFunction_eq_zero_re};
\item \code{LFunction_eq_zero_re_of_isPrimitive}.
\end{itemize}
They cover all positive moduli and both primitive and imprimitive characters, with the principal-character pole treated separately.\cite{repo,repo_readme}

The repository specifies Lean 4.35.0 and Mathlib revision
\begin{center}\code{56a005855ea17e0b0958f73ec801ae140b0a1679}.\end{center}
Its README and the independent verification page report compilation, axiom inspection, kernel replay and comparison of the public statements. The reported axioms are \code{propext}, \code{Classical.choice}, and \code{Quot.sound}.\cite{repo_readme,verification}

We inspected the pinned source and public types. Lean and Lake were unavailable on the execution path, and we did not perform a new full build or independent kernel replay. Thus statements about those earlier checks are attributed to their published records. The new proofs in this report have not been checked by Lean.

The endpoint does not quantify over arbitrary Hecke characters of arbitrary fields, arbitrary Artin representations, or general automorphic $L$-functions. An abelian factorization can transfer Dirichlet GRH to certain Dedekind zeta functions, but a nonabelian factorization contains genuinely additional functions. Renaming an input character does not supply a new zero theorem.

\subsection{Three arithmetic estimates and a one-sided remainder}

The reading guide isolates three inputs: short integral forms, enough independent forms, and a lower bound for integral forms. Their comparison is the mechanism to transfer.\cite{guide}

\begin{enumerate}
\item \textbf{Divisibility makes forms arithmetically short.} Derivatives of theta functions carry prime-power divisibility. Chains of theta lifts repeat the gain. After normalization, many source vectors give integral vectors whose archimedean norm bounds contain the weighted prime sum.
\item \textbf{A second moment supplies rank.} Spanning estimates ensure that the short vectors occupy almost all of the target space. Without independence, many small vectors say little about a determinant.
\item \textbf{Congruences and twists control denominators.} The lower bound uses determinants of lifts, heights, and integral structures shared across auxiliary characters. Selmer estimates are used to find characters that keep congruence losses within budget.
\item \textbf{The product formula compares the two sides.} The same determinant is constrained by its archimedean size and finite-place divisibility. The principal terms cancel, leaving an upper bound for a centered prime remainder.
\item \textbf{A one-sided analytic criterion excludes zeros.} A smoothed prime sum has a meromorphic transform whose poles encode zeros. Positivity and a suitable auxiliary quadratic twist prevent the relevant residues from cancelling.
\end{enumerate}

Schematically, the remainder at an identity residue class is
\[
 R_{q,d}(x)=
 \sum_{\substack{n\le x\\ n\equiv1\pmod q}}
 \Lambda(n)(1-\chi_{-d}(n))
       \left(\frac1n-\frac1x\right)
 -\frac{\log x}{\varphi(q)}-c_{q,d}.
\]
The actual parameters obey special congruence and CM-field conditions. The guide gives a chain length $m$ between $\sqrt x$ and $2\sqrt x$, and an eventual upper bound of size $x^{-1/2+6\varepsilon}$ after the remaining parameters are chosen. The analytic argument must work with the required family of auxiliary primes $d$, so that one with a zero-free prescribed rectangle can be selected. A bound for a single convenient twist does not automatically replace this quantifier.\cite{guide}

\subsection{Which statements are reusable?}

The abstract theorem \code{New.TwistedLattice.covolume_add_index_nonneg} has the form
\begin{equation}\label{eq:covolume-source}
 0\le c+2\delta+n'(b_{\rm short}+b_{\rm denom})+b_{\rm volume}.
\end{equation}
Its parameters describe the short family, an index bound, a completion to a basis, and a height bound for a ratio of determinants. This is an actual abstract transfer interface. The concrete \code{integral_forms_long} and \code{star_inequality} instantiate it with the source's special cyclotomic-CM data, not with arbitrary $L$-functions.\cite{covolume_source}

The auxiliary twisting statements also carry substantial hypotheses. The readable development's general theorem requires residual absolute irreducibility, finiteness of conjugate-dual congruence terms, finite local determinant terms, and a scalar inertia condition of the form
\[
 \rho(\gamma)=\rho(\gamma^{-c})=\zeta I,
 \qquad \zeta^2-1\text{ a unit}.
\]
The specialized main-summand predicate asks for a stronger family of nontrivial scalar powers. These conditions cannot be suppressed when testing a Tate twist or an elliptic-curve representation. The compact development packages its own conditions using balanced data and auxiliary twists.\cite{selmer_source}

One further condition is central below. The predicate \code{IsFromU1} forces triviality on the \emph{whole decomposition group} at specified places, not just on inertia. At a prime outside the conductor, that means its Frobenius must also map to one. The source proves this explicitly as \code{IsFromU1.eq_one_of_mem_decompGroup}.\cite{decomp_source}

Finally, the source's uniformly bounded odd-cocycle module is a custom object. The comparison with the usual strict Selmer or Tate--Shafarevich group is additional mathematics. A bound for that custom module alone does not establish a new vanishing theorem for an algebraic $K$-group.\cite{odd_source}

\section{When local torsion data become global twists}\label{sec:characters}

A torsion-avoidance lemma chooses roots of unity satisfying polynomial inequalities. A global twisting theorem needs those roots to be values of an arithmetic character. There are three separate quantities: the orders of the local inertia values, the order of the global character, and its conductor. The following elementary theorem identifies their first compatibility obstruction.

\subsection{The finite-group theorem}

\begin{theorem}[Exact extension criterion]\label{thm:extension}
Let $N\ge1$, and let
\[
 0\longrightarrow A\longrightarrow B\longrightarrow C\longrightarrow0
\]
be an exact sequence of finite abelian $p$-groups, with $p^aC=0$. Let $D$ be a subgroup of $B$, and let $\lambda:A\to\mu_{p^N}$ be a character.
\begin{enumerate}
\item An extension $\chi:B\to\mu_{p^{N+a}}$ trivial on $D$ exists if and only if
\begin{equation}\label{eq:compat-raised}
              \lambda(A\cap D)=1.
\end{equation}
\item An extension $\chi:B\to\mu_{p^N}$ trivial on $D$ exists if and only if
\begin{equation}\label{eq:compat-exact}
              \lambda\bigl(A\cap(D+p^NB)\bigr)=1.
\end{equation}
\end{enumerate}
Here $\lambda(H)=1$ means that $\lambda$ is trivial on every element of $H$.
\end{theorem}

\begin{proof}
A character from a subgroup of a finite abelian $p$-group extends to $\mu_{p^\infty}$. An elementary construction adjoins one generator $b$ at a time: if the least positive integer $r$ with $rb$ in the old subgroup is $r=p^j$, choose an $r$th root of the already assigned value at $rb$. This gives a well-defined extension and terminates.

Condition~\eqref{eq:compat-raised} is necessary for an extension killing $D$. If it holds, $\lambda$ descends to $\bar A=(A+D)/D\subseteq B/D$. Extend it into $\mu_{p^\infty}$. The quotient $(B/D)/\bar A$ is a quotient of $C$, hence is killed by $p^a$. For each $\bar b\in B/D$,
\[
 \chi(\bar b)^{p^{N+a}}
 =\lambda(p^a\bar b)^{p^N}=1.
\]
Thus the extension takes values in the required finite group.

For the second assertion, a character of exponent dividing $p^N$ kills $p^NB$ as well as $D$, giving necessity. Conversely, under \eqref{eq:compat-exact}, $\lambda$ descends to the image of $A$ in $B/(D+p^NB)$. Extend it as above. This quotient is killed by $p^N$, so its character values automatically lie in $\mu_{p^N}$.
\end{proof}

The permitted increase in order is bounded by the exponent of the quotient, not by its cardinality. When $D=0$, the proportion of local $\mu_{p^N}$-valued characters that extend with that same exponent is exactly
\begin{equation}\label{eq:proportion}
 \frac{1}{[A\cap p^NB:p^NA]}.
\end{equation}
Indeed, all such local characters form the dual of $A/p^NA$, while the extendible characters form the dual of $A/(A\cap p^NB)$. This is a useful correction to any argument that samples the unrestricted local torsion grid and then assumes that the same sampling measure applies to global characters.

\subsection{The arithmetic exact sequence}

Let $E$ be a CM field, with complex conjugation $c$. Fix an odd prime $p$ and a squarefree modulus
\[
 \mathfrak m=\prod_{j=1}^t\mathfrak v_jc(\mathfrak v_j),
\]
where all the primes are distinct, $\mathfrak v_j\ne c(\mathfrak v_j)$, and none lies above $p$. Put $q_j=\N\mathfrak v_j$. For a finite abelian $p$-group with $c$-action, the minus part is the image of the idempotent $(1-c)/2$.

\begin{theorem}[CM ray realization]\label{thm:ray}
Set
\[
 R=\prod_{j=1}^t\F_{q_j}^{\times}[p^\infty],\qquad
 B=\Cl_{\mathfrak m}(E)[p^\infty]^{-},\qquad
 C=\Cl(E)[p^\infty]^{-}.
\]
Let $W$ be the image of $\mu(E)[p^\infty]$ in the minus local residue group $R$. Then $W$ is cyclic, and there is an exact sequence
\begin{equation}\label{eq:ray}
 0\longrightarrow R/W\longrightarrow B\longrightarrow C\longrightarrow0.
\end{equation}
Consequently Theorem~\ref{thm:extension} applies with $A=R/W$. A proposed local tuple must first annihilate $W$. Prescribed splitting conditions are imposed by the appropriate subgroup $D$ of ray classes.
\end{theorem}

\begin{proof}
Start with the usual ray class exact sequence
\[
 \mathcal O_E^\times\longrightarrow
 (\mathcal O_E/\mathfrak m)^\times\longrightarrow
 \Cl_{\mathfrak m}(E)\longrightarrow\Cl(E)\longrightarrow0.
\]
The Chinese remainder theorem identifies each conjugate residue pair with two copies of $\F_{q_j}^\times$, exchanged by $c$. Its minus $p$-part consists of $(u,u^{-1})$ and is one copy of $\F_{q_j}^\times[p^\infty]$. Primary decomposition and the minus idempotent are exact direct-summand operations on the finite groups.

The minus part of the unit image needs care. For any unit $u$, the quotient $u/c(u)$ is an algebraic integer unit all of whose conjugates have absolute value one. Kronecker's theorem makes it a root of unity. On the $p$-part, taking its square root is an automorphism, since $p$ is odd. Thus the minus unit image is contained in the image of $\mu(E)[p^\infty]$. The reverse inclusion holds because every such root is a unit and $c$ acts on it by inversion. The unit image is therefore exactly $W$, proving \eqref{eq:ray}.
\end{proof}

This proof uses the standard ideal-theoretic class-field correspondence to interpret characters of $B$ as finite anticyclotomic Hecke characters. The foundational exact sequence is classical; the point here is to retain all its terms when transferring the torsion-selection argument.\cite{milne}

If $p\nmid|\mu(E)|$, then $W=0$ and the local variables are unrestricted before the class and splitting conditions are imposed. If $p\nmid h_E$, then additionally $C=0$, so the unrestricted minus local tuple determines a unique $p$-primary ray character. Neither simplification can be assumed in general. In particular, the upstream field contains a primitive root of order $5040$; its $p=3,5,7$ cases require the $W$ term.

\subsection{Conductor and powers}

An extension supplied by Theorem~\ref{thm:ray} has conductor dividing $\mathfrak m$. At each prime of this squarefree modulus, a nontrivial local residue character has conductor exponent one. Therefore, if all prescribed local characters are nontrivial,
\begin{equation}\label{eq:conductor}
 \mathfrak f(\chi)=\mathfrak m,\qquad
 \N\mathfrak f(\chi)=\prod_{j=1}^t q_j^2.
\end{equation}
Increasing the allowed global order does not force an increase in this conductor.

If the local character at pair $j$ has exact order $p^{n_j}$, the first $M$ positive powers of $\chi$ all retain the full conductor if and only if $p^{n_j}>M$ for every $j$. This assertion concerns local orders, even when the global order is larger.

These conclusions still do not construct a small Selmer twist. They solve the character-extension part of that construction. Tied cocycle relations, nonvanishing polynomial witnesses, and valuation budgets are additional inputs.

\subsection{How many independent parameters survive?}

\begin{proposition}[A fixed-exponent parameter bound]\label{prop:grid}
Suppose $N\ge1$ and $R=\prod_{j=1}^t C_{p^{h_j}}$ with $h_j\ge N$. Let $C$ require $r$ generators, let $W$ require $w$ generators, and let the required-triviality subgroup $D\subseteq B$ have a chosen generating set of size $s$. Then the image of
\[
 \Hom(B/D,\mu_{p^N})\longrightarrow\Hom(R,\mu_{p^N})
\]
contains a free $\Z/p^N\Z$-submodule of rank at least
\[
                     \max(0,t-w-r-s).
\]
For the CM unit subgroup, $w$ is zero or one.
\end{proposition}

\begin{proof}
Choose invariant-factor generators $c_i$ of $C$, of orders $p^{a_i}$, and lifts $b_i$ in $B$. Write each relation $p^{a_i}b_i$ as an element of $R/W$, and choose a representative in $R$. Assign the value zero to every $b_i$, writing characters additively in $\Z/p^N\Z$. Assigning values to the $t$ local generators then gives a character precisely when the unit relations and these $r$ relations are respected. Requiring the chosen $s$ elements of $D$ to vanish adds at most $s$ further linear equations. Thus the kernel of a matrix with at most $w+r+s$ rows over $\Z/p^N\Z$ maps into the restriction image. Smith normal form leaves at least the asserted number of free columns.
\end{proof}

The bound $t-r$ in the case $W=D=0$ is sharp. For $a\ge1$, take
\[
 B=C_{p^{N+a}}^r\times C_{p^N}^{\,t-r},\qquad
 A=p^aC_{p^{N+a}}^r\times C_{p^N}^{\,t-r}.
\]
The first $r$ restriction coordinates lie in $p^a(\Z/p^N\Z)$ and have smaller exponent; the last $t-r$ give all free coordinates of exponent $p^N$.

The proposition does not say that every original inertia coordinate can vary independently, or even that every one is nonzero somewhere on the displayed submodule. A polynomial which is nonzero on the full local parameter space may vanish identically on the compatible subgroup. Applying the upstream torsion-avoidance result after this restriction requires a fresh witness estimate. This is one of the precise remaining bridges.

\section{An exact conductor minimum over \texorpdfstring{$\Q(\sqrt{-23})$}{Q(sqrt(-23))}}\label{sec:example}

We now compute a small example completely. It demonstrates both a genuine extension obstruction and the cost of retaining a source condition.

Let
\[
 E=\Q(\sqrt{-23}),\qquad
 \theta=\frac{1+\sqrt{-23}}2,\qquad
 \theta^2-\theta+6=0.
\]
The reduced primitive forms of discriminant $-23$ are
\[
 (1,1,6),\quad(2,-1,3),\quad(2,1,3).
\]
The reduced-form inequality $a\le\sqrt{23/3}<3$ makes this list exhaustive. Hence $\Cl(E)\simeq C_3$, entirely in the minus part. The only roots of unity are $\pm1$.

The prime $\mathfrak q=(3,\theta-1)$ is nonprincipal. Indeed, the norm form is $u^2+uv+6v^2$; for $v\ne0$ it is at least $23v^2/4>3$, and for $v=0$ it is a square. Set
\[
 \beta=1+2\theta=2+\sqrt{-23},\qquad \N\beta=27.
\]
Its reduction is zero at $\mathfrak q$ and nonzero at the conjugate prime above $3$, so
\begin{equation}\label{eq:q3-cube}
                       \mathfrak q^3=(\beta).
\end{equation}
In particular the class of $\mathfrak q$ generates $\Cl(E)$.

\subsection{Why conductor 13 fails}

At $13$, the roots of $T^2-T+6$ are $5$ and $9$, so $13$ splits. The minus local $3$-group has order three. The ray exact sequence gives a group $B$ of order nine:
\[
 0\longrightarrow C_3\longrightarrow B\longrightarrow C_3\longrightarrow0.
\]
At the root $\theta=5$, the residue of $\beta/c(\beta)$ is $4$ modulo $13$, and
\[
                           4^4\equiv9\not\equiv1\pmod {13}.
\]
Thus its $3$-primary component is nontrivial. By \eqref{eq:q3-cube}, the triple of a minus lift of the class above $3$ is a nonzero local inertia element. The lift therefore has order nine and $B\simeq C_9$.

Its local inertia subgroup is $3C_9$. Every character $C_9\to\mu_3$ kills that subgroup. A prescribed nontrivial local cubic character therefore needs global order nine, even though its conductor norm remains $13^2=169$. This proves that the order increase in Theorem~\ref{thm:extension} can be necessary in an actual number field.

The splitting condition above $3$ is more restrictive. The minus Frobenius at $3$ generates $C_9$; requiring its value to be one kills the whole ray group. Thus \emph{no nontrivial character of any $3$-power order} on this fixed support satisfies that condition. Enlarging the allowed character order does not repair it.

\subsection{The first support which works}

\begin{theorem}[Exact minimum]\label{thm:minimum}
Among nontrivial finite $3$-power anticyclotomic characters of $E$ which are trivial on the decomposition groups above $3$ and ramify only at primes lying in split rational conjugate pairs, the minimum conductor norm is
\[
                              151^2=22801.
\]
Exactly two characters attain this minimum. Both have order three and conductor $(151)$.
\end{theorem}

\begin{proof}
For a split rational prime $q\ne3$, the minus local $3$-group is cyclic of order $3^{h_q}$, where $h_q=v_3(q-1)$. If $h_q=0$, no such character can ramify at $q$. Assume henceforth $h_q\ge1$.

Let $B_q$ be the minus $3$-ray group modulo $(q)$, and let $b$ be a minus lift of Frobenius above $3$. Its image generates the ordinary class group. Consequently $B_q$ is generated by its local group $R_q$ and $b$, with the single additional class relation
\[
                         3b=r_q\in R_q.
\]
Up to inversion and multiplication by a $3$-adic unit, $r_q$ is the $3$-primary residue of $\beta/c(\beta)$, by \eqref{eq:q3-cube}. Killing Frobenius therefore gives
\begin{equation}\label{eq:ray-quotient}
                 B_q/\langle b\rangle
                    \simeq R_q/\langle r_q\rangle.
\end{equation}
It is nontrivial exactly when $\beta/c(\beta)$ is a cube modulo $q$. Choosing the conjugate root in the residue field only inverts the ratio and does not change this test.

The exact finite calculation below includes every split prime $q\le151$ with $q\equiv1\pmod3$. The root $t$ is a root of $T^2-T+6$ and
\[
        u_q=\frac{1+2t}{3-2t}\pmod q,\qquad
        z_q=u_q^{(q-1)/3}\pmod q.
\]
\begin{center}
\begin{tabular}{rrrrrr}
\toprule
$q$ & $t$ & $u_q$ & $h_q$ & $z_q$ & $|B_q/\langle b\rangle|$\\
\midrule
13  &5 &4  &1&9 &1\\
31  &8 &13 &1&5 &1\\
73  &30&54 &2&64&1\\
127 &28&78 &2&19&1\\
139 &17&102&1&42&1\\
151 &43&79 &1&1 &3\\
\bottomrule
\end{tabular}
\end{center}
At $151$ the quotient is $C_3$, giving exactly two nontrivial characters. Their inertia is nontrivial, so the conductor is exactly $(151)$.

For minimality, unramified characters are killed by the condition at $3$ because the class of $\mathfrak q$ generates $\Cl(E)$. A character with one rational prime in its support and smaller conductor would occur in an earlier row and has just been excluded. A character with at least two distinct rational support primes has their product at least $13\cdot31=403>151$. Anticyclotomicity makes both conjugate primes ramify together. At residue characteristic $q\ne3$, a finite $3$-power character is trivial on the higher unit groups, which are pro-$q$; every ramified conductor exponent is one. Thus a multi-prime support has conductor norm at least $403^2>151^2$. This proves the claim without a squarefree-conductor assumption being added separately.
\end{proof}

The certificate uses only exact reduced-form enumeration, norms, modular roots, and modular exponentiation. The class-field interpretation is proved above, not delegated to floating-point computation. The file \code{checks/check_ray_characters.py} produces the full table and a JSON certificate.

\subsection{What this example repairs}

A two-prime alternative also illustrates the role of the support. At $13$ and $31$, both local groups are $C_3$ and both components of the relation $3b$ are nonzero. After changing their generators, killing $b$ imposes the relation $(1,1)$ on $C_3^2$. The two nontrivial compatible characters have local exponent vectors $(1,-1)$ and $(-1,1)$. Their conductor norm is $(13\cdot31)^2=162409$. The one-prime support $151$ is cheaper.

These constructions satisfy the elementary anticyclotomic, tame-ramification, and decomposition-group conditions. They do not verify the full upstream auxiliary-twist predicate. For example, for a primitive cubic value,
\[
                    v_3(1-\zeta_3)=\frac12.
\]
Summing over the two primes of a single split pair already gives one; over the two pairs above $13$ and $31$ it gives two. An arbitrarily small prescribed valuation budget is not met. Selmer and special-value conditions remain separate as well.

The example is a completed small arithmetic optimization problem and a rigorous test case for the transfer interface. It is not presented as a new solution of a broad conjecture or as a literature-priority claim.

\section{An explicit conductor bound and its limits}\label{sec:conductor}

The source's torsion and prime-selection arguments control the number and character-theoretic properties of auxiliary primes. They do not give a norm bound for primes satisfying every required Frobenius and Selmer condition. One can nevertheless give a completely explicit bound for the simpler, unrestricted realization problem.

\begin{theorem}[Bounded conductor for an abelian CM base]\label{thm:bounded-conductor}
Let $E/\Q$ be an abelian CM extension contained in $\Q(\zeta_f)$. Let $p$ be odd, $N\ge1$, and $t\ge1$. Write
\[
 |\mu(E)[p^\infty]|=p^{w_0},\qquad
 p^a\Cl(E)[p^\infty]^{-}=0.
\]
Let $S$ be a finite set of rational primes to avoid. Set
\[
 M=\operatorname{lcm}\left(f,p^{N+w_0},\prod_{q\in S}q\right).
\]
Choose a prime $\ell>t$ not dividing $M$, and put
\[
 Q=M\ell,\qquad P=(\varphi(Q)\log Q)^2.
\]
Assuming Dirichlet GRH, there exist $t$ distinct rational primes $q_j\le P$, splitting completely in $E$ and avoiding $S\cup\{p\}$, with the following property. Choose one conjugate pair above each $q_j$. Every tuple of local characters of order dividing $p^N$ on those pairs extends to an anticyclotomic character $\chi$ with
\[
 \ord(\chi)\mid p^{N+a},\qquad
 \N\mathfrak f(\chi)\le P^{2t}.
\]
If all the selected local characters have exact order $p^N$, then their conductor is exactly the selected modulus, and its first $m$ powers have that same conductor whenever $m<p^N$.
\end{theorem}

\begin{proof}
Choose $t$ distinct nonzero residues modulo $\ell$. For each, the Chinese remainder theorem gives a reduced residue class modulo $Q$ which is $1$ modulo $M$ and has the chosen value modulo $\ell$. The least-prime bound of Lamzouri, Li and Soundararajan gives a prime in each class at most $P$, since $Q>3$.\cite{lls} They are distinct because the residues modulo $\ell$ are distinct. Their congruence modulo $f$ ensures complete splitting in $E$; their congruence modulo the other factors ensures avoidance and
\[
                       v_p(q_j-1)\ge N+w_0.
\]

For the corresponding local $p$-group $R$, every coordinate of the root-unit image $W$ has order at most $p^{w_0}$. In a cyclic factor of order $p^{h_j}$, with $h_j\ge N+w_0$, such elements lie in $p^{h_j-w_0}R_j\subseteq p^NR_j$. Therefore
\[
                              W\subseteq p^NR.
\]
Every local character into $\mu_{p^N}$ automatically annihilates $W$. Apply Theorems~\ref{thm:extension} and \ref{thm:ray} with $D=0$. The conductor and power assertions follow from \eqref{eq:conductor}.
\end{proof}

This strengthens a construction which simply excludes primes dividing the roots of unity of $E$: deeper local residue groups absorb that unit relation. Among supports with $h_j\ge N$, the extra depth $h_j\ge N+w_0$ is also necessary for \emph{every independent} local $\mu_{p^N}$ tuple to kill the unit image. Reduction away from $p$ injects $\mu(E)[p^\infty]$ into each residue field, so the projection of $W$ is exactly the order-$p^{w_0}$ subgroup; it lies in $p^NR_j$ precisely at the displayed depth. A tuple supported on one coordinate detects any failure.

The construction also supplies arbitrarily many independent local coordinates using distinct rational primes, without an exponentially growing recursive exclusion modulus. It is a synthesis of a known explicit prime bound with the arithmetic extension theorem, not a new least-prime theorem.

There is an elementary lower restriction. A tame inertia character of order $p^N$ requires $p^N\mid q_j-1$, hence $q_j\ge p^N+1$. A nontrivial conjugate pair therefore has conductor norm at least $(p^N+1)^2$. No conductor bound independent of the requested local order can hold.

\subsection{The part of the small-valuation condition this supplies}

Normalize $v_p(p)=1$. At a primitive $p^N$th root,
\[
 v_p(1-\zeta_{p^N})=\frac{1}{(p-1)p^{N-1}}.
\]
For $1\le j\le m<p^N$,
\[
 v_p(1-\zeta_{p^N}^{\,j})
     =\frac{p^{v_p(j)}}{(p-1)p^{N-1}}.
\]
Thus a tuple of $t$ conjugate pairs whose local values all have exact order $p^N$ satisfies the uniform tame bound
\[
 \sum_{\mathfrak v\mid\mathfrak m}
 v_p(1-\chi(\tau_{\mathfrak v})^j)
 \le
 \frac{2t\,p^{\lfloor\log_p m\rfloor}}{(p-1)p^{N-1}},
 \qquad 1\le j\le m.
\]
Here $\tau_{\mathfrak v}$ is a tame generator, and primitivity makes this valuation independent of its choice. Choosing $N$ large makes this particular cost small, with the conductor bound from Theorem~\ref{thm:bounded-conductor}. This recovers one useful quantitative component of the source's twisting problem.

\subsection{The condition this construction does not supply}

Theorem~\ref{thm:bounded-conductor} sets $D=0$. It does not impose splitting at the primes above $p$. The obstruction in Section~\ref{sec:example} shows that this omission can be decisive even at tiny discriminant. Nor do the congruences on $q_j$ enforce the Frobenius neighborhoods governing the source's tied Selmer relations. Those can involve nonabelian extensions.

One may quotient by the required subgroup $D$ and apply the exact criterion, or use Proposition~\ref{prop:grid} to exhibit a surviving family. Either route still needs a proof that the determinant or special-value witness does not vanish on that family. Selecting a prime in a rational progression alone does not prove this.

Prescribed tame ramification through relations among Frobenius elements is a developed subject, including the Gras--Munnier criterion and its modern proofs.\cite{hmr} Effective Chebotarev results can be considered once the governing field, degree and discriminant are actually bounded.\cite{thorner} They cannot be replaced silently by the Dirichlet endpoint of the present Lean proof.

For a nonabelian base field, even finding small completely split primes requires additional analysis. For $p=2$, the minus idempotent used here is unavailable; the unit and conjugation analysis must be replaced. Neither extension is claimed in this report.

\section{Deeper confluent certificates and exact sparse families}\label{sec:confluent}

The incoming proof of the last quintic case uses scalar congruences and divided differences to make a finite test valid for arbitrarily large exponents. Its stated sufficient condition is $p>k+1$, where $k$ is the total number of powers. Retaining a factorial numerator gives the following stronger statement.

\begin{theorem}[Higher-depth, blockwise certificate]\label{thm:confluent}
Let $p$ be any prime, $s,m\ge1$, and let $A$ be a finite free commutative $\Z_{(p)}$-algebra. Suppose
\[
                    x^m=a+p^sD,\qquad p\nmid a,
\]
with $x,D\in A$ and $a\in\Z$. Let $\mathcal E$ be a finite set of distinct nonnegative integers, and let $n_r$ count its elements congruent to $r$ modulo $m$, for $0\le r<m$. Assume
\begin{equation}\label{eq:block-bound}
                     n_r<p^s\qquad\text{for every }r.
\end{equation}
If the columns
\[
                x^rD^j,\qquad 0\le r<m,\quad0\le j<n_r,
\]
are linearly independent in $A/pA$ over $\F_p$, then the powers $x^e$, $e\in\mathcal E$, are linearly independent in $A\otimes\Q$ over $\Q$.
\end{theorem}

\begin{proof}
Within a residue block write the exponents as $e_i=r+mt_i$. Divide the associated columns by $a^{t_i}$. They become
\[
                         x^r(1+p^sD/a)^{t_i}.
\]
For $j=0,\ldots,n_r-1$, take the Newton divided difference of order $j$ at the first $j+1$ distinct integers $t_i$, multiplied by $j!/p^{sj}$. These are invertible triangular column operations over $\Q$.

The coefficient of $(D/a)^i$ in the new column is
\[
 j!p^{s(i-j)}
       [t_0,\ldots,t_j]\binom{T}{i}.
\]
It is zero for $i<j$ and one for $i=j$. For $i>j$, its valuation is at least
\begin{equation}\label{eq:factorial-bound}
                       s(i-j)-v_p(i!/j!).
\end{equation}
Indeed, $i!\binom{T}{i}$ is an integer polynomial, and divided differences of integer polynomials at integer arguments are integers. For a monomial this is the complete homogeneous symmetric-polynomial identity, and linearity gives the general case.

Let $S_p(n)$ denote the base-$p$ digit sum. Since $j\le n_r-1\le p^s-2$,
\[
 S_p(j)\le s(p-1)-1,\qquad S_p(i)\ge1.
\]
Set $d=i-j\ge1$ and $C=s(p-1)-1\ge0$. Legendre's formula gives
\begin{align*}
 (p-1)\bigl(sd-v_p(i!/j!)\bigr)
 &=Cd-S_p(j)+S_p(i)\\
 &\ge C(d-1)+1>0.
\end{align*}
Every later coefficient is therefore divisible by $p$. The transformed column is integral and reduces to $x^r(D/a)^j$. Its factor $a^{-j}$ is a nonzero scalar modulo $p$.

A rational dependence between integral columns in a finite free $\Z_{(p)}$-module can be scaled to have integral coefficients with at least one unit. Reduction would contradict the assumed independence. Reverse the invertible column operations to obtain the result.
\end{proof}

The proof uses finite binomial expansions; no $p$-adic exponential or logarithm is required. Thus $p=2$ causes no convergence problem. The bound is on each residue block, not on total support. At depth one it needs $n_r<p$, already weaker than the preceding total-support condition.

\subsection{Sharpness of the strict threshold}

The strict inequality in \eqref{eq:block-bound} cannot be replaced universally by a weak inequality. Put $k=p^s$ and
\[
 A=\Z_{(p)}[D]/(D^k+D^{k-1}-1),\qquad x=1+kD,\quad m=1.
\]
The $k$ confluent columns $1,D,\ldots,D^{k-1}$ form a basis modulo $p$. Nevertheless
\[
 F(X)=(X-1)^k+k(X-1)^{k-1}-k^k
\]
vanishes at $x$ and has exactly $k$ nonzero coefficients. The coefficient of $X^{k-1}$ is zero; all the coefficients at exponents $0,1,\ldots,k-2,k$ are nonzero. At $p=3,s=1$, the example becomes the irreducible relation $x^3-3x-25=0$ after $D^3+D^2-1=0$.

This proves sharpness in the theorem's finite-free-algebra generality. It does not assert that the chosen order equals $\Z_{(p)}[x]$; in fact that narrower equality need not hold.

\subsection{Two exact seven-term families}

\begin{theorem}[Powered singular moduli]\label{thm:powered}
Let $\alpha$ be any singular modulus of discriminant $-87$, and let $\beta$ be any singular modulus of discriminant $-152$. Then, for every integer $t\ge1$,
\[
                     s(\alpha^{63t})=7,\qquad
                     s(\beta^{182t})=7.
\]
The exponents and rational coefficients of the annihilating polynomials are unrestricted.
\end{theorem}

\begin{proof}
Both ring class polynomials have degree six. Their exact integer coefficients and the certificates below are in \code{checks/confluent_certificates.json}. Their identification as class polynomials is inherited from the prior report's interval-certified CM calculations, whose complete verifier was rerun in this audit.\cite{prior_sparse}

In $A=\Z_{(2)}[X]/(H_{-87})$, exact reduction gives
\[
 x^{63}\equiv 1+16(x+x^2+x^3)\pmod {32A}.
\]
Thus $x^{63}=1+16D$ with $D\bmod2=x+x^2+x^3$. The six columns $1,D,\ldots,D^5$, represented with entries $0,1$ in the power basis, have integer determinant $-1$. They are independent modulo $2$, and $6<2^4$. Apply Theorem~\ref{thm:confluent} to any six distinct exponents divisible by $63$.

Likewise, in $\Z_{(3)}[X]/(H_{-152})$,
\[
 x^{182}\equiv26+9(2x^2+x^3+2x^5)\pmod {27A}.
\]
For $D=(x^{182}-26)/9$, the six columns $1,D,\ldots,D^5$ have the following representative matrix modulo $3$:
\[
 \begin{pmatrix}
 1&0&2&2&0&0\\
 0&0&2&0&2&2\\
 0&2&2&1&1&2\\
 0&1&0&0&0&0\\
 0&0&0&0&2&1\\
 0&2&1&2&0&1
 \end{pmatrix}.
\]
Its integer determinant is $-2$, hence is nonzero modulo $3$, and $6<3^2$.

For either powered number, a rational relation with at most six terms would select at most six such powers. A smaller dependent set could be enlarged to six distinct exponents, contradicting the established independence. On the other hand, the number lies in a six-dimensional $\Q$-algebra, so seven powers are dependent. This gives support exactly seven. Replacing $63$ or $182$ by any positive multiple retains the argument.
\end{proof}

These are exact infinite-exponent statements, not bounded searches. The modular powers were verified both by polynomial multiplication and by an independently written basis-shift recurrence. The matrices and determinants were checked with integer arithmetic.

The result does not prove $s(-87)=7$ or $s(-152)=7$. A six-term relation for the original singular modulus could involve several different exponent residue classes. That remaining possibility is precisely what a complete residue-cover search would have to exclude.

\subsection{Why this is a useful next step}

The local idea fits the source's arithmetic theme: retain divisibility which a coarser size estimate discards, then certify independence through a determinant. Here the implementation is elementary and the proofs are short enough to formalize independently. It gives a concrete partial result in the next support problem, while avoiding any claim that an analogous theta-lift or Selmer construction has already been built for a new $L$-function.

The deeper certificate should be included in a future search over six-term exponent patterns, especially at small primes. The output of that search must be an exhaustive finite cover of all residue patterns, including repeated residues; checking large exponent ranges alone would not prove the unbounded statement.

\section{A different way to pay for missing dimensions}\label{sec:rank}

The source makes the short family span almost the whole target space. The abstract determinant theorem reveals what that hypothesis pays for: completing the family to a basis. It is useful to retain that cost explicitly when designing a new arithmetic application.

\begin{proposition}[Rank and completion cost]\label{prop:rank}
Assume the hypotheses of the upstream corrected-covolume inequality \eqref{eq:covolume-source}. Let the target dimension be $d=n+n'$, with $n>0$, and let $m>0$. Put $C=b_{\rm short}+b_{\rm denom}$. Suppose
\[
 c\le n(A-mR+u),\qquad b_{\rm volume}\le-dA+B.
\]
Then
\begin{equation}\label{eq:rank-cost}
 \boxed{\displaystyle
 mR\le u+\frac{2\delta+B}{n}
             +\frac{d-n}{n}(C-A).}
\end{equation}
No additional sign assumptions on $A,C,R,u,B$ are needed.
\end{proposition}

\begin{proof}
Substitute the two upper bounds into \eqref{eq:covolume-source}:
\[
 0\le -nmR+nu+2\delta+B+(d-n)(C-A).
\]
Rearrange and divide by $n>0$.
\end{proof}

This elementary rearrangement is useful because the deficit is charged by $C-A$, the \emph{excess} completion cost, rather than by $C$ alone. Write $\varrho=n/d$. For an error scale $E_0>0$, suppose
\[
 \varrho\ge\rho>0,\quad
 u\le UE_0,\quad \delta/d\le IE_0,\quad B/d\le JE_0,\quad
 (1-\varrho)(C-A)_+\le KE_0.
\]
Then
\begin{equation}\label{eq:rank-error}
 R\le\frac{U+(2I+J+K)/\rho}{m}\,E_0.
\end{equation}
There are consequently two ways to achieve the needed accuracy: make the missing dimension fraction extremely small, or prove that a completion has excess cost $O(E_0)$. In the latter case a fixed positive fraction of independent target forms can suffice. If $C\le A$, the deficit term is favorable. A uniform lower bound on $n/d$, or a compensating improvement in the other errors, is still needed.

\subsection{An explicit weakening inside the existing proof}

In the source's notation, let
\[
 d=\mathrm{targetDim},\quad A=\mathrm{mainAmount},\quad
 u=fa(\log b)^4,\quad C=2x_1^2b^2+b.
\]
Before using its near-full-rank hypothesis, the proof of \code{integral_forms_long} has already obtained
\begin{equation}\label{eq:source-before-rank}
 0\le c+2\delta+(d-n)C+d(b+u-A).
\end{equation}
The completing vectors come from a globally spanning auxiliary family, with squared-norm exponent $2x_1^2b^2$, scalar denominator-height bound $b$, and the displayed volume-height bound. These arithmetic inputs must be retained; an arbitrary completion is insufficient.\cite{long_source}

Equation~\eqref{eq:source-before-rank} gives
\[
 c+2\delta\ge d(A-b-u)-(d-n)C.
\]
Thus the source's conclusion
\[
                         c+2\delta\ge d(A-3u)
\]
already follows under the completion budget
\begin{equation}\label{eq:weaker-spanning}
                  (d-n)C\le d(2u-b).
\end{equation}
Since the source has $b\le u$, a convenient sufficient condition is $(d-n)C\le du$. The original argument instead obtains the smaller allowance $(d-n)C\le db$ from the loss bound $d-n\le4d/b^2$ and its scale inequalities.

Keeping those same completion and height estimates, it is therefore sufficient to replace that particular loss bound by
\[
 \frac{d-n}{d}\le
      \min\left\{\frac12,\frac{u}{2x_1^2b^2+b}\right\}.
\]
The first term ensures $n\ge d/2$ when the other errors are divided by $n$. With $A\ge0$, $\delta\le dI_0$, $b\le u$, and the short-family upper bound from Proposition~\ref{prop:rank}, this gives
\[
                             mR\le7u+4I_0.
\]
Indeed, use $B=d(b+u)$ in \eqref{eq:rank-cost}, drop the favorable term $-(d-n)A$, and bound the numerator by $d(2I_0+3u)$.

This is a genuine weakening of a sufficient hypothesis in the written source argument, not an improvement to the already proved RH endpoint. For the same expensive completion the missing fraction still tends to zero. For a new target, however, the calculation suggests a distinct research task: construct cheaper integral completing vectors and control their twist scalars, instead of attempting immediately to reproduce the full upstream spanning estimate.

\subsection{What a nonabelian transfer would still need}

The abstract inequality requires injective twist maps, a family and a completion that together form a basis, integral vectors after a common scalar, norm bounds at all complex conjugates, and height bounds at the allowed finite places. The determinant ratio is tied to this same lattice. An abstract frame, a representation-theoretic identity, or a numerical eigenvalue bound does not supply these arithmetic properties.

With the source-like error scale $E_0=fa(\log b)^4$ and $m\asymp\sqrt x$, \eqref{eq:rank-error} has normalized size $O((\log b)^4/\sqrt x)$. Converting this into the final prime remainder still needs the separate comparison of the arithmetic gain with the relevant Frobenius statistic. That is precisely the missing step in the earlier nonabelian criteria.

\section{Which open questions should be pursued next?}

\subsection{First priority: the six-term singular-modulus problem}

The prior classification \eqref{eq:prior-classification} gives a concrete next question:
\[
 \text{If }h(\Delta)\ge6,\quad\text{must }s(\Delta)\ge7?
\]
The two exact powered families prove partial cases of the corresponding support-seven phenomenon. They do not settle the original six-term question.

An effective next stage would start with the remaining low-degree orders and certify all exponent residue patterns using the stronger blockwise test. The required output is finite and precise: exact class-polynomial coefficients, scalar congruences modulo $p^{s+1}$, and a covering argument showing that every candidate six-exponent pattern is rejected by at least one rank test. A successful global statement would additionally need a large-discriminant reduction, as the five-term report supplied. The finite and infinite parts should be proved separately and then joined.

This target has three advantages. It is already connected to a published question about rational sparse equations for singular moduli;\cite{blp} it has exact arithmetic certificates; and the new local theorem changes which small-prime tests are available. The generic fewnomial literature provides context, while the CM conjugates and the finite residue covering supply the special arithmetic content.\cite{fewnomials}

\subsection{Second priority: an arithmetic family of constrained twists}

For the Selmer route, the next useful deliverable is not another isolated torsion point. It is a family that simultaneously has:

\begin{enumerate}
\item an explicit ray-group presentation and required-triviality subgroup $D$;
\item a proved nonvanishing or bounded-valuation witness on the compatible character subgroup;
\item controlled local orders so the first required powers remain ramified;
\item a quantitative support-prime construction meeting the actual Frobenius conditions; and
\item a bound for the associated conductor and coefficient-field ramification.
\end{enumerate}

Theorems~\ref{thm:extension}--\ref{thm:bounded-conductor} solve parts of this interface and show exactly why the parts do not collapse into one assertion. The $151$ example is a complete elementary benchmark. A harder benchmark would replace its single class relation by the tied cocycle relations of one explicitly chosen Selmer datum.

The auxiliary primes in the source are selected in prescribed Frobenius neighborhoods. One must identify the governing extension before deciding which zero theorem or effective Chebotarev bound is available. Rational congruences only control abelian information. This is where a purported conductor bound based only on Dirichlet GRH can fail.

\subsection{Nonabelian GRH: aim at one arithmetic inequality}

Earlier reports already reduce several nonabelian zeta questions to explicit one-sided prime-sum estimates. The cubic field $\Q(\alpha)$ with $\alpha^3-\alpha-1=0$, and its $S_3$ splitting field, are natural small benchmarks; a class-sensitive Hecke character over $\Q(\sqrt{-23})$ occurs there. The new character calculation concerns auxiliary ray characters of that quadratic field. It does not prove a zero theorem for its class-character $L$-function.

Further progress should be measured by a new arithmetic input: an integral lattice with the right Frobenius-sensitive divisibility, an independent family, a cheap completion, or a determinant lower bound with a controlled error. Proposition~\ref{prop:rank} allows these inputs to be budgeted before an ambitious construction is attempted.

The repository endpoint cannot simply be applied to the nonabelian factor. Nor does finding anticyclotomic characters imply that the theta lifts or Selmer hypotheses needed for its determinant comparison have been constructed.

\subsection{\texorpdfstring{$K$}{K}-groups, Vandiver, and BSD}

For the $K$-theory route, the most valuable next result is the comparison between the source's custom odd-cocycle module and the conventional cohomological object, with all local conditions and exceptional primes recorded. After that, a finite Hecke determinant certificate may become a proof of a genuine arithmetic statement. Before it, the computation is a candidate test.

For Vandiver-type questions, parity, field variation, and the distinction between a normalized length and actual vanishing remain material. Passing to a more ramified coefficient field can make normalized lengths small without killing a module. The global-order increase of Section~\ref{sec:characters} should therefore be recorded explicitly, not hidden in a change of coefficient ring.

For BSD, the general twisting theorem is not a theorem about arbitrary elliptic-curve Selmer groups. The scalar-inertia and conjugate-dual conditions must be established for the actual representation. Anticyclotomic twisting can preserve a conjugate-dual obstruction rather than remove it. A direct substitution remains unsupported.

\subsection{Research verdict}

There is real transferable mathematics here. The completed CM classification and the new powered families show a tractable direction in which arithmetic divisibility and determinants give exact results beyond familiar RH consequences. The character calculation gives another concrete output and corrects a tempting but false identification between local torsion data and global twists of the same order.

The path to a new broad zero theorem remains open at its arithmetic core. The practical way forward is to prove a missing construction or estimate for one chosen target, with a numerical error budget such as \eqref{eq:rank-cost}, and to stop counting equivalent analytic criteria as resolutions.

\section{Evidence, reproducibility, and formalization}

\subsection{Pinned inputs}

The repository snapshots are:
\begin{itemize}
\item \code{lsernau/grh-lean}:\\
\code{72fead102f663d5c7a407964a13dc0cf9d977207}.
\item \code{VladimirReshetnikov/ProveIt}:\\
\code{48efe71a726dddc5a505534edaa4c0fd190416b9}.
\end{itemize}
The directory inventory retains the incoming archive names, original Git blob identifiers and sizes. The two previously existing CM archives are included unchanged under \code{prior/}, together with their fresh replay logs. A separate source manifest records the selected upstream files and SHA-256 hashes. The package manifest covers all deliverable files.

\subsection{What was actually checked}

\begin{center}
\small
\begin{tabular}{p{.34\linewidth}p{.57\linewidth}}
\toprule
Item & Verification in this report\\
\midrule
Public Lean theorem scope & Pinned statement and source inspection; earlier external build reports attributed\\
Full upstream Lean proof & No new kernel replay; no local build claimed\\
Older four-term CM report & Complete supplied reproduction script rerun successfully\\
Current five-term CM report & Complete reproduction script rerun, including two last-quintic and independent sextic implementations\\
New ray-character minimum & Exact norm, form, finite-group and residue certificates; independently checked calculations\\
New confluent instances & Exact class-polynomial modular arithmetic with two recurrences, rank and integer determinants\\
General new lemmas & Written proofs with internal review; not Lean-verified\\
Literature priority & Not established; known ingredients attributed and new synthesis distinguished\\
\bottomrule
\end{tabular}
\end{center}

No finite numerical experiment is used to infer an all-exponent statement. The finite rank tests imply those statements through Theorem~\ref{thm:confluent}. Likewise, the minimum conductor calculation uses the proved ray-group quotient and an exhaustive finite list below the proposed minimum, plus a lower bound for all multi-prime supports.

\subsection{Running the package}

From the extracted archive directory, run
\begin{verbatim}
python3 reproduce.py
\end{verbatim}
This uses only the Python standard library for the new finite certificates. The generated JSON and text records are under \code{evidence/}. To replay the previous CM packages as well, use
\begin{verbatim}
python3 reproduce.py --include-prior
\end{verbatim}
The older verifiers also need their documented dependencies, including a C++17 compiler for the exhaustive residue sieve. Their original archives and reproduction files are retained rather than silently rewritten.

Rebuild this report with
\begin{verbatim}
latexmk -pdf -interaction=nonstopmode -halt-on-error report.tex
\end{verbatim}
The distributed PDF was compiled from the included source and rendered for layout review.

\subsection{A dependency-ordered Lean plan}

The smallest useful formalization starts with the finite-group extension criterion, keeping the constrained subgroup $D$ explicit. Next come the finite-free-algebra divided-difference lemma, the factorial-valuation estimate, and the transfer from a modular rank certificate to rational independence. These can be stated independently of the upstream automorphic machinery.

The CM ray sequence requires units, conjugation, finite primary components and global reciprocity. Its treatment of $p\mid|\mu(E)|$ should be formalized at the same time as the unrestricted case. The exact $-23$ calculation can then be attached through a small certified class-group presentation.

For the sparse examples, the modular certificates are small, but a fully Lean-certified singular-modulus theorem additionally needs the class-polynomial identities. Reusing their Python interval certificates is not itself a Lean proof. The dependency map under \code{formalization/} makes this boundary explicit.

The covolume rearrangement can be formalized as a short ordered-field lemma once the existing theorem is imported. Proving its arithmetic inputs for a new target remains the much larger task.

\clearpage
\begin{thebibliography}{99}\small\raggedright
\bibitem{repo}
L. Sernau and collaborators, \emph{GRH for Dirichlet $L$-functions, in Lean 4},
revision \code{72fead102f663d5c7a407964a13dc0cf9d977207}, 25 November 2026.
\href{https://github.com/lsernau/grh-lean/tree/72fead102f663d5c7a407964a13dc0cf9d977207}{Repository snapshot}.
Public types: \href{https://github.com/lsernau/grh-lean/blob/72fead102f663d5c7a407964a13dc0cf9d977207/grh/src/Statements.lean}{Statements.lean};
exports: \href{https://github.com/lsernau/grh-lean/blob/72fead102f663d5c7a407964a13dc0cf9d977207/grh/src/Main.lean}{Main.lean}.
\bibitem{repo_readme}
The same repository, \href{https://github.com/lsernau/grh-lean/blob/72fead102f663d5c7a407964a13dc0cf9d977207/grh/README.md}{readable edition README} and
\href{https://github.com/lsernau/grh-lean/blob/72fead102f663d5c7a407964a13dc0cf9d977207/grh-compact/README.md}{compact edition README}.
\bibitem{verification}
A. V. Sutherland, \emph{Verification files for the shorter Lean proof of RH and GRH},
published check records for the pinned revision.
\url{https://math.mit.edu/~drew/shortgrh/}.
These are external verification records, not a new replay by this report.
\bibitem{guide}
Sernau and collaborators, \emph{Reading guide}, pinned revision.
\href{https://github.com/lsernau/grh-lean/blob/72fead102f663d5c7a407964a13dc0cf9d977207/grh/GUIDE.md}{GUIDE.md}.
For the analytic criterion see
\href{https://github.com/lsernau/grh-lean/blob/72fead102f663d5c7a407964a13dc0cf9d977207/grh/src/New/Reduction/Zeros/Criterion.lean}{Criterion.lean}.
\bibitem{covolume_source}
Sernau and collaborators,
\href{https://github.com/lsernau/grh-lean/blob/72fead102f663d5c7a407964a13dc0cf9d977207/grh/src/New/Reduction/Lattices/TwistedLattice.lean}{TwistedLattice.lean},
especially \code{covolume_add_index_nonneg}.
\bibitem{long_source}
Sernau and collaborators,
\href{https://github.com/lsernau/grh-lean/blob/72fead102f663d5c7a407964a13dc0cf9d977207/grh/src/New/Determinant/Twist/LongForms.lean}{LongForms.lean},
especially \code{mul_completion_cost_le} and \code{integral_forms_long}.
\bibitem{selmer_source}
Sernau and collaborators, compact edition
\href{https://github.com/lsernau/grh-lean/blob/72fead102f663d5c7a407964a13dc0cf9d977207/grh-compact/src/New/Selmer/TwistTheorem.lean}{TwistTheorem.lean},
\href{https://github.com/lsernau/grh-lean/blob/72fead102f663d5c7a407964a13dc0cf9d977207/grh-compact/src/New/Selmer/RootsOfUnityPoints.lean}{RootsOfUnityPoints.lean}, and
\href{https://github.com/lsernau/grh-lean/blob/72fead102f663d5c7a407964a13dc0cf9d977207/grh-compact/src/New/Selmer/TrianglePrimes.lean}{TrianglePrimes.lean}.
For the readable-edition hypotheses see
\href{https://github.com/lsernau/grh-lean/blob/72fead102f663d5c7a407964a13dc0cf9d977207/grh/src/New/Congruences/Selmer/GeneralTheorem.lean}{GeneralTheorem.lean} and
\href{https://github.com/lsernau/grh-lean/blob/72fead102f663d5c7a407964a13dc0cf9d977207/grh/src/New/Congruences/Selmer/SelmerData.lean}{SelmerData.lean}.
\bibitem{decomp_source}
Sernau and collaborators, compact edition
\href{https://github.com/lsernau/grh-lean/blob/72fead102f663d5c7a407964a13dc0cf9d977207/grh-compact/src/Lib/ClassField/Places/SelmerDefs.lean}{SelmerDefs.lean},
\href{https://github.com/lsernau/grh-lean/blob/72fead102f663d5c7a407964a13dc0cf9d977207/grh-compact/src/Lib/ClassField/Places/Conjugation.lean}{Conjugation.lean}, and
\href{https://github.com/lsernau/grh-lean/blob/72fead102f663d5c7a407964a13dc0cf9d977207/grh-compact/src/New/Selmer/Corepresentable.lean}{Corepresentable.lean}.
\bibitem{odd_source}
Sernau and collaborators,
\href{https://github.com/lsernau/grh-lean/blob/72fead102f663d5c7a407964a13dc0cf9d977207/grh/src/Lib/Cyclotomic/OddSha.lean}{OddSha.lean},
especially the exact definition and scope of \code{exists_length_shaOdd_le}.
\bibitem{prior_sparse}
V. Reshetnikov, ProveIt incoming archive,
\emph{Five-term equations for singular moduli: an exact classification and lower bounds in arbitrary degree},
\href{https://github.com/VladimirReshetnikov/ProveIt/blob/48efe71a726dddc5a505534edaa4c0fd190416b9/docs/incoming/grh_sparse_singular_moduli.zip}{grh\_sparse\_singular\_moduli.zip}.
Original Git blob \code{0d5a83ba283e3d100168c9569cb981c010bafe7c}.
The prior package, source and fresh replay log accompany this report.
\bibitem{prior_four}
V. Reshetnikov, ProveIt incoming archive,
\emph{Sparse equations for singular moduli and the sign of weighted prime sums},
\href{https://github.com/VladimirReshetnikov/ProveIt/blob/48efe71a726dddc5a505534edaa4c0fd190416b9/docs/incoming/grh_cm_four_terms_and_prime_bias.zip}{grh\_cm\_four\_terms\_and\_prime\_bias.zip}.
Original Git blob \code{3f043cce34fd9dda2512e0dd013637029e3537de}.
\bibitem{blp}
Y. Bilu, F. Luca and A. Pizarro-Madariaga,
\emph{Trinomials, singular moduli and Riffaut's conjecture},
Ann. Sc. Norm. Super. Pisa Cl. Sci. (5) 23 (2022), 2003--2048.
\href{https://arxiv.org/abs/2003.06547}{arXiv:2003.06547}.
\bibitem{fewnomials}
M. Avenda\~no and T. Krick,
\emph{Sharp bounds for the number of roots of univariate fewnomials},
J. Number Theory 131 (2011), 1209--1228.
\href{https://arxiv.org/abs/1008.4808}{arXiv:1008.4808}.
\bibitem{milne}
J. S. Milne, \emph{Class Field Theory}, version 4.03 (2020),
Chapter V, ray class groups and ideal-theoretic reciprocity.
\url{https://jmilne.org/math/CourseNotes/CFT.pdf}.
\bibitem{lls}
Y. Lamzouri, X. Li and K. Soundararajan,
\emph{Conditional bounds for the least quadratic non-residue and related problems},
especially Corollary 1.2.
\href{https://arxiv.org/abs/1309.3595}{arXiv:1309.3595}, version 3.
The bound used is $P(a,q)\le(\varphi(q)\log q)^2$ for $q>3$ under Dirichlet GRH.
\bibitem{hmr}
F. Hajir, C. Maire and R. Ramakrishna,
\emph{On tame $\Z/p\Z$ extensions with prescribed ramification}.
\href{https://arxiv.org/abs/2208.05007}{arXiv:2208.05007}.
\bibitem{thorner}
J. Thorner and A. Zaman,
\emph{An explicit bound for the least prime ideal in the Chebotarev density theorem},
Algebra Number Theory 11 (2017), 1135--1197.
\href{https://arxiv.org/abs/1604.01750}{arXiv:1604.01750}.
\end{thebibliography}

\end{document}
